\section{Fourier multiplication on \texorpdfstring{$\Zpstar{p}$}{(Z/p)*}}
\label{sec:mod-mul-algorithm}

We trained a one-layer transformer with the architecture and hyperparameters of
Nanda et al.\ \cite{nanda2023progress} on the task $(a, b) \mapsto a \cdot b
\bmod p$ for $p = 113$.  The network groks.  We claim that the learned network
implements the algorithm in Section~\ref{subsec:mul-algorithm}, an analog of
Nanda's Fourier multiplication algorithm composed with the discrete-log
isomorphism of Lemma~\ref{lem:dlog-iso}.

\subsection{Setup}
\label{subsec:setup}

The network has one transformer block, $d_{\text{model}} = 128$, $4$ attention
heads of dimension $d_{\text{head}} = 32$, MLP width $d_{\text{mlp}} = 512$,
ReLU activation, no LayerNorm.  Inputs are tokenized as $[a, b, =]$; logits are
read at position $2$.  Training: full batch, AdamW with $\eta = 10^{-3}$,
$\lambda = 1.0$, $30\%$ of the $p^2$ pairs as the training set.  Grokking onset
at epoch~$\approx 3500$; convergence by epoch~$\approx 5000$.

\subsection{The discrete-log isomorphism}
\label{subsec:dlog}

\begin{lemma}[Discrete-log isomorphism]
\label{lem:dlog-iso}
Let $p$ be an odd prime and $g$ a primitive root modulo $p$.  The map
\begin{equation}
  \dlog_g \colon \Zpstar{p} \;\xrightarrow{\;\sim\;}\; \Zp{p-1},
  \qquad g^{k} \longmapsto k,
  \label{eq:dlog-iso}
\end{equation}
is a group isomorphism, under which multiplication on $\Zpstar{p}$ becomes
addition on $\Zp{p-1}$:
\begin{equation}
  \dlog_g(a \cdot b) \;\equiv\; \dlog_g(a) + \dlog_g(b) \pmod{p-1}.
  \label{eq:mult-to-add}
\end{equation}
\end{lemma}

\noindent The Pontryagin dual of $\Zp{p-1}$ is the cyclic group of multiplicative
characters $\chi_k \colon \Zp{p-1} \to \C^\times$, $\chi_k(j) = e^{2\pi i k j /
(p-1)}$, indexed by $k \in \Zp{p-1}$.  Pulling these back through $\dlog_g$
yields the multiplicative characters of $\Zpstar{p}$:
\begin{equation}
  \chi_k(a) \;=\; \exp\!\left(\tfrac{2\pi i\, k\, \dlog_g(a)}{p-1}\right),
  \qquad
  \chi_k(a b) = \chi_k(a)\,\chi_k(b).
  \label{eq:mult-character}
\end{equation}

\subsection{The Fourier multiplication algorithm on \texorpdfstring{$\Zpstar{p}$}{(Z/p)*}}
\label{subsec:mul-algorithm}

We claim that the network implements the following algorithm.  Let
$\omega_k := 2\pi k / (p-1)$ and fix a small set $\mathcal{K} \subset \{1,
\ldots, \lfloor (p-1)/2 \rfloor\}$ of \emph{key frequencies}
(empirically, $|\mathcal{K}| \approx 4$).

\begin{enumerate}[leftmargin=2em]
  \item \emph{Embedding.}  For each input token $a$ and each $k \in
        \mathcal{K}$, write
        \begin{equation}
          \WE[\,\cdot\,, a] \;\supseteq\;
            \cos(\omega_k\, \dlog_g a), \;
            \sin(\omega_k\, \dlog_g a)
          \quad\text{(as two distinct $d_{\text{model}}$-directions).}
          \label{eq:embed-claim}
        \end{equation}
  \item \emph{Attention + MLP.}  Combine the embeddings of positions $0$ and
        $1$ and apply the angle-addition identities
        \begin{align}
          \cos(\omega_k(\dlog_g a + \dlog_g b))
          &= \cos(\omega_k \dlog_g a)\cos(\omega_k \dlog_g b)
            - \sin(\omega_k \dlog_g a)\sin(\omega_k \dlog_g b),
            \label{eq:trig-cos} \\
          \sin(\omega_k(\dlog_g a + \dlog_g b))
          &= \sin(\omega_k \dlog_g a)\cos(\omega_k \dlog_g b)
            + \cos(\omega_k \dlog_g a)\sin(\omega_k \dlog_g b).
            \label{eq:trig-sin}
        \end{align}
        By \eqref{eq:mult-to-add}, the left-hand sides equal
        $\cos(\omega_k\, \dlog_g(ab))$ and $\sin(\omega_k\, \dlog_g(ab))$.
  \item \emph{Unembed.}  For each candidate output $c$, the logit is
        \begin{equation}
          \mathrm{logit}(c)
          \;\approx\;
          \sum_{k \in \mathcal{K}}
            \cos\!\bigl(\omega_k\, (\dlog_g(ab) - \dlog_g c)\bigr).
          \label{eq:logit-formula}
        \end{equation}
        Each summand is maximized at $\dlog_g c \equiv \dlog_g(ab) \pmod{p-1}$,
        i.e.\ at $c \equiv a b \pmod p$; constructive interference selects the
        correct $c$.
\end{enumerate}

We refer to this as the \emph{Fourier multiplication algorithm} on $\Zpstar{p}$.
It is Nanda's modular-addition algorithm in the log domain, applied to the
isomorphic group $\Zp{p-1}$.

\subsection{Where each step is implemented}
\label{subsec:implementation}

\begin{itemize}[leftmargin=*]
  \item \eqref{eq:embed-claim} is realized by the embedding matrix $\WE \in
        \R^{d_{\text{model}} \times p}$.  Empirically, projecting $\WE$ onto
        the multiplicative character basis \eqref{eq:mult-character} concentrates
        $\geq 97\%$ of its $\ell_2$-energy on the $2|\mathcal{K}|$ rows
        corresponding to $\{\cos(\omega_k \cdot), \sin(\omega_k \cdot) :
        k \in \mathcal{K}\}$.
  \item \eqref{eq:trig-cos}--\eqref{eq:trig-sin} are realized jointly by the
        attention layer (which copies the per-position character values to the
        final position) and the MLP, which computes the bilinear products via
        $\mathrm{ReLU}$ followed by $\Wout$.  Each MLP neuron's centered
        activation on the $(a, b)$ grid is sparse in the 2D character basis,
        concentrated on the four bigrams corresponding to a single $k \in
        \mathcal{K}$.
  \item \eqref{eq:logit-formula} is realized by $\WU \Wout$, whose rows are
        well-approximated by linear combinations of $\cos(\omega_k \dlog_g c)$,
        $\sin(\omega_k \dlog_g c)$ for the same $k \in \mathcal{K}$.
\end{itemize}

\subsection{Relation to prior work}
\label{subsec:relation}

Nanda et al.\ \cite{nanda2023progress} established this algorithm for modular
\emph{addition} on $\Zp{p}$.  Our result is the same algorithm composed with
the discrete-log isomorphism \eqref{eq:dlog-iso}; equivalently, it is Fourier
multiplication on $\Zp{p-1}$ pulled back to $\Zpstar{p}$.  Gromov
\cite{gromov2023grokking} derives an analytic closed form for the weights of a
2-layer MLP grokking modular addition; the analogous derivation for our
multiplicative setting amounts to the same expressions reindexed by $\dlog_g$.
Mallinar et al.\ \cite{mallinar2024emergence} show that recursive feature
machines learn block-circulant features for modular addition.  In our setting,
the analogous statement is that the model's weight matrices are block-circulant
\emph{after reindexing by $\dlog_g$}; we leave the explicit check of this
statement to future work.
